\documentclass[%
  twocolumn
%   , hidelinks
%   , hidempi
%   , linenumbers
]{mpi2015-cscpreprint}

%%%%%%%%%%%%%%%%%%%%%%%%%%%%%%%%%%%%%%%%%%%%%%%%%%%%%%%%%%%%%%%%%%%%%%%%%%%%%%%%
% INDIVIDUAL PACKAGES.                                                         %
%%%%%%%%%%%%%%%%%%%%%%%%%%%%%%%%%%%%%%%%%%%%%%%%%%%%%%%%%%%%%%%%%%%%%%%%%%%%%%%%

% Proper hyphenation.
\usepackage[american]{babel}

% Graphical packages.
\usepackage{graphicx}
\usepackage{bm}
% Math packages.
\usepackage{amssymb}
\usepackage{amsthm}


%%%%%%%%%%%%%%%%%%%%%%%%%%%%%%%%%%%%%%%%%%%%%%%%%%%%%%%%%%%%%%%%%%%%%%%%%%%%%%%%
% MAIN DOCUMENT.                                                               %
%%%%%%%%%%%%%%%%%%%%%%%%%%%%%%%%%%%%%%%%%%%%%%%%%%%%%%%%%%%%%%%%%%%%%%%%%%%%%%%%


\begin{document}
  
%%%%%%%%%%%%%%%%%%%%%%%%%%%%%%%%%%%%%%%%%%%%%%%%%%%%%%%%%%%%%%%%%%%%%%%%%%%%%%%%
% PAPER INFORMATION.                                                           %
%%%%%%%%%%%%%%%%%%%%%%%%%%%%%%%%%%%%%%%%%%%%%%%%%%%%%%%%%%%%%%%%%%%%%%%%%%%%%%%%

\title{Internal Report: Inference of Numerical Schemes}
  
\author[$\ast$]{Süleyman Yıldız}
\affil[$\ast$]{Address of first author.\authorcr
  \email{yildiz@mpi-magdeburg.mpg.de}, \orcid{0000-0000-0000-0000}}
  
\author[$\dagger$]{Leonidas Gkimisis}
\affil[$\dagger$]{Address of second author.\authorcr
  \email{gkimisis@mpi-magdeburg.mpg.de}, \orcid{0000-0000-0000-0000}}

\author[$\dagger\dagger$]{Peter Benner}
\affil[$\dagger\dagger$]{Address of second author.\authorcr
  \email{benner@mpi-magdeburg.mpg.de}, \orcid{0000-0000-0000-0000}}
  
\shorttitle{Example short title}
\shortauthor{F. Author, S. Author}
\shortdate{}
  
\keywords{keyword1, keyword2, keyword3}

\msc{MSC1, MSC2, MSC3}
  
\abstract{%
  This is the preprint template of the Department for Computational Methods
  in Systems and Control Theory (CSC) of the Max Planck Institute for Dynamics
  of Complex Technical Systems, Magdeburg.
  This template is supposed to be the default for all preprints with main or
  corresponding authors at our institute.
  In the following, the use of the template is explained in more details
  with some example for easy use.}

% \novelty{This assesses what the new and significant contribution of this
%   research is.}

\maketitle


%%%%%%%%%%%%%%%%%%%%%%%%%%%%%%%%%%%%%%%%%%%%%%%%%%%%%%%%%%%%%%%%%%%%%%%%%%%%%%%%
% PAPER CONTENT.                                                               %
%%%%%%%%%%%%%%%%%%%%%%%%%%%%%%%%%%%%%%%%%%%%%%%%%%%%%%%%%%%%%%%%%%%%%%%%%%%%%%%%

\subsection{Least Squares problem formulation}

\textcolor{red}{To be removed}

Focusing on a grid node $i$, we denote the data on that node as $\mathbf{u}$ and the set of ``neighbouring" nodes as $q(i)$. We also consider a continuous time formulation and a linear inference problem. Then, the local inference of the discrete operators requires solving the following least-squares problem:

\begin{equation}
\min _{\Vec{a}_i}{\left\|\Vec{a}_i^T \mathbf{u}_{q(i)}-\frac{du}{dt}|i\right\|_{2}},
\end{equation}


It is evident that one least-squares problem should be solved for every grid node $i$ and thus the above problem scales with the original full-order-model dimension. However it is valuable to examine the extension of the above problem to systems with a more general polynomial structure. As an example, the incompressible Navier-Stokes equations have a quadratic structure [+++]. Polynomial nonlinearities are also used to approximate general nonlinear systems [+++]. The following generalizes also to systems with control terms (e.g. bilinear structure).



Then, we need to locally infer all involved discretized operators, by solving the following least-squares problem:

\begin{equation} 
\label{LS_noreg}
\min _{\Vec{a}_i}{\left\|\Vec{a}_i^T \mathcal{D}_i-\mathbf{u}_i^{k+1}\right\|_{2}},
\end{equation}

where matrix $\mathcal{D}_i$ contains data coming from the adjacent nodes of node $i$, for the different polynomial terms:

\begin{equation}
\label{Dmat}
\mathcal{D}_i^T=
\begin{bmatrix}
  \mathbf{u_{q(i)}} \\
  \mathbf{u}_{q(i)}^2\\
  \mathbf{u}_{q(i)}^3\\
  ...\\
\end{bmatrix}
.
\end{equation}


\subsection{Computational cost}

It is interesting to examine the computational cost involved. For each least-squares problem with $n_t$ timesteps and $n_x$ discretized operator coefficients, the least-squares solution cost asymptotically scales as $\mathcal{O}(n_x^2 \times n_t)$. We assume that in system +++ is overdetermined $(n_t>n_x)$. Considering $s$ adjacent nodes for node $i$, then the non-repeating adjacent node products for a polynomial term of order $k$, will be $s \choose k$. This leads to the conclusion that $n_x$ scales with $s^k$. 

Thus the overall cost is $\mathcal{O}(s^{2k} \times n_t)$. Then locally identifying the discrete operators for all nodes would scale as $\mathcal{O}(n \times s^{2k} \times n_t)$. We observe that the argument of adjacency-based sparsity allows scaling proportionally to $n$. The polynomial degree $k$ depends on the structure of the system we aim to infer, while the number of adjacent nodes $s$  to be considered would also vary depending on the system at hand. However, one could argue that for the same spatial dimension (e.g. for two-dimensional problems), $s$ will not vary significantly.


Eventually, the solution of +++ is an inferred numerical scheme which aims to mimic the action of the corresponding discretized operator to $u_i$. In an idealized scenario, one could thus think that locally inferring the operators at every point $i$ is redundant. Indeed, if we limit ourselves to uniform grids, then we expect that the solution of +++ will be the same for all internal points in the grid. As mentioned in [+++], the boundary conditions should be treated independently. In such cases, it seems that the involved computational cost could be significantly reduced and no longer scale with $n$. Finally, it should be mentioned that the $n$ involved least-squares problems can be solved independently and are thus fully parallelizable. Thus, the corresponding offline clock time can be significantly reduced.


Examining the data-driven inference of discretized operators might be also specifically interesting for model reduction. The high offline cost involved is leveraged by several advantages that could assist/complement projection-based methods. In particular, working on the full-order scale allows for more robust models [+++] that inherit intrusive ROM properties. Notably, there is flexibility in boundary condition treatment, since Dirichlet boundary conditions can be directly enforced to the corresponding nodes. Moreover, considering a model with a known parametric dependence, local inference could potentially assist in handling bifurcation phenomena. If the data-driven model structure is in line with the corresponding discretized PDE, the non-smooth parametric dependencies of the system dynamics could be fitted to the inferred operators, while a projection basis could be chosen at a later stage. Finally, performing the step of inference on the full-order model bypasses issues of non-Markovian dynamics, which can often cause instabilities for non-intrusive ROMs [+++].

In the following, we focus on one-dimensional linear and nonlinear problems on a uniform grid. We thus examine in detail the mathematical properties of the solution of (+++). First, we provide with theoretical results on linear-advection-dominated problems which give valuable insight on the local numerical scheme inference. We then proceed to investigate numerical simulations for several linear and nonlinear academic examples which highlight different aspects of the local inference problem. The relation of space/time discretization, stability and consistency constraints enforcement, ill-posedness of the least-squares problem and the effect of data quality are extensively discussed and specific solutions are proposed.


\section{A toy example: Advection equation}
Focusing on symmetric schemes ($l=k$), the main error originates from $\langle u_o,\frac{du}{dt}|_o \rangle$. We observe that if the system at hand is energy-preserving, then $\langle u_o,\frac{du}{dt}|_o \rangle=0$. In that case, the approximate solutions \ref{approxD} and \ref{approxC} are consistent. 

If we consider a linear advection problem, $\frac{du}{dt}=\gamma \frac{du}{dx}$, then $\beta_m$ coincides with the 1st-order central differences scheme.

On the other side, we can consider a diffusion problem with one slow and one fast scale on a domain $x \in [0,1]$, such that the solution is pde solution is:

\begin{equation}
    u(x,t)=e^{-\pi^2 t}sin(\pi x) + b e^{-\pi^2 q^2 t} sin (q \pi x)
\end{equation}

%Lets show again the orthogonality of u with dudx

We choose $q=4z$ with $z \in \mathcal{N}$ and derive the solution $\beta_m$ for $x_c=1/2$. One can see that this point is such that $\langle \frac{du}{dx}|_o,\frac{du}{dt}|_o \rangle \approx 0$, whereas  $\langle u_o,\frac{du}{dt}|_o \rangle=-\pi^2 a$. Then the obtained symmetric scheme would be $\beta_m \approx -\pi^2/3$.

From this example, we aim to highlight the importance of the sampling point $x_c$ selection. It is evident that this point should be such that $u_o$ and $\frac{du}{dt}|_o$ are almost orthogonal. 
\subsection{Theoretical Insights}

In order to get some analytical insight on the data-driven inference of a numerical scheme, we focus on solutions originating from one-dimensional (1D) advection-dominated problems. In particular, we assume the data is originating from a solution of the type

\begin{equation}
\label{refeq}
g(x,t)=sin(x+ct) f(x)
\end{equation}

for $x$ in $[-L,L]$. We specify that $f(x)$ is either an odd or an even function (or that $f^2(x)$ is an even function). We discretize $x$ with $\Delta x$ and focus on the data of one spatial point $x_c$. Assuming a symmetric stencil of $2n+1$ points and using a typical trigonometric identity, we write:

\begin{multline}
\label{decomp}
g(x_c+k\Delta x,t)=sin(x_c+ct)cos(k\Delta x) f(x)+\\
+cos(x_c+ct)sin(k\Delta x) f(x)
\end{multline}

where $k=[-n_x,...,n_x]$. The terms in \ref{decomp} are orthogonal functions. Thus, it actually corresponds to the analytical analog of singular value decomposition for a data matrix of $g(x_c+k\Delta x,t)$, with $n_t$ rows and $2n+1$ columns. This property allows us to analytically examine the data-driven solution for the inferred, symmetric numerical scheme in this case.

In particular, since the scheme is symmetric, the first SVD mode will correspond to the solution of the midpoint, (thus $sin(x_c+ct)$). 
%(talk about statistical stationarity)

The basis functions can be written out as follows:


\begin{align}
\Phi_1(t)=\frac{sin(x_c+ct)}{|| sin(x_c+ct) || } ,\\
\Phi_2(t)=\frac{cos(x_c+ct)}{|| cos(x_c+ct) || } 
\end{align}

Similarly, for the $x$ dimension:

% not correct when f(x) plays an important role.

\begin{align}
\Psi_1(k)=\frac{cos(k \Delta x) f(x_c+k \Delta x)}{|| cos(k \Delta x) f(x_c+k \Delta x) || } ,\\
\Psi_2(k)=\frac{sin(k \Delta x) f(x_c+k \Delta x)}{|| sin(k \Delta x) f(x_c+k \Delta x) || }
\end{align}

Then the singular values are:

\begin{align}
\sigma_1=|| sin(x_c+ct) || || cos(k \Delta x) f(x_c+k \Delta x) || ,\\
\sigma_2=|| cos(x_c+ct) || || sin(k \Delta x) f(x_c+k \Delta x) ||
\end{align}

+++how accurate are these estimates? Add maybe the following figure: L2 error on sigmas, w.r.t T and dissipation (surf).


Then, any data matrix $\mathcal{D} \in 		\mathbb{R} ^{n_t \times {2n_x+1}}$  with snapshots from the solution of \ref{refeq} can be decomposed into

\begin{equation}
    \mathcal{D}=[\Phi_1 \Phi_2]
    \begin{bmatrix}
\sigma_1 & 0 \\
0 & \sigma_2
\end{bmatrix}
    \begin{bmatrix}
\Psi_1 \\
\Psi_2
\end{bmatrix}
\end{equation}

We then aim to solve the following least-squares problem:

\begin{equation}
\label{lsadv}
    \mathcal{D} \Vec{v}= \Vec{b}
\end{equation}

where in continuous time $\Vec{b}=\dot{g}(x_c,t)$ and in discrete time $\Vec{b}=g(x_c,t+\Delta t)$.



\subsubsection{Continuous time}

Taking the derivative of \ref{refeq}, it can be shown that

\begin{equation}
\label{deriv}
\dot{g}(x_c,t)=\Phi_2(t)\underbrace{c ||cos(x_c+ct)|| f(x_c)}_{\mbox{\alpha(x_c)}} 
\end{equation}

Substituting into \ref{lsadv}, the analytical solution of the least-squares problem can be obtained:

\begin{equation}
    \Vec{v}=\sigma_2^{-1} {\Psi_2}^T \alpha(x_c)
\end{equation}

We observe that the second singular value scales the solution, while the second right singular vector $\Psi_2$ is giving the solution to the least squares problem. 

The independence of the solution with respect to the first singular value in cases of continuous time can be generalized for any statistically stationary dataset. In particular, for such systems, it can be shown that the first singular vector $\sigma_1 \Phi_1$ is approximately the dataset average [++++]. Then if the solution is written as $u(t)=\tilde{u}(t)+\bar{u}$, with $\bar{u}$ being the average, the derivative will be $\dot{u}(t)=\dot{\tilde{u}}(t)$. One can then see that \ref{deriv} will result to a least-squares solution $\Vec{v}$ with a linear combination of $\Psi_2 ... \Psi_n$. Hence, $\Vec{v}$ will be independent of the average.

Focusing on datasets coming from \ref{refeq}, we substitute the corresponding expressions to finally get:

\begin{equation}
    v(k)=\frac{c f(x_c)f(k)sin(k \Delta x)}{{|| sin(i \Delta x) f^2(i)||}^2}
\end{equation}


In the case of the simple advection equation with sinusoidal initial data, $f(x)=1$. Considering $n=1$, we want to compare the obtained data-driven scheme to the central difference scheme.

\begin{align}
v(1)={}& -v(-1)=\frac{c}{2 sin(\Delta x)} \\
v(0)={}& 0
\end{align}

This indicates that the inferred scheme is always conservative ($\sum_i{v(i)}=0$), irrespective of the choice of $\Delta x$. However, it will be $0$-order convergent. Examining the error with respect to the central difference scheme coefficient, we get that

\begin{equation}
v(1)-\frac{c}{2 \Delta x}=-\frac{c \Delta x}{6 + 6\mathcal{O}(\Delta x ^2)}.
\end{equation}

This showcases a linear convergence of the inferred scheme to the true, 1st order central difference scheme. As a result, the inferred scheme converges to a second order accurate scheme, only at the limit of $\Delta x \to 0$. 

++++ one simulation plot to show the linear decrease. \textcolor{red}{Suleyman}

\subsubsection{Discrete time}
\label{disc_time}

The case of the discrete time advection equation with sinusoidal initial data can provide with further insight on the solution properties of \ref{lsadv}. Indeed, in practically any application, the analytical derivative information is not available.

In this case,  $\Vec{b}=g(x_c,t+\Delta t)$, which can be written as:

\begin{equation}
g(x_c,t+\Delta t)=[\Phi_1 \Phi_2]     
\begin{bmatrix}
|| sin(x_c+ct) || \; cos(c \Delta t) \\
|| cos(x_c+ct) || \; sin(c \Delta t)
\end{bmatrix}
f(x_c)
\end{equation}

We denote 

\begin{equation}
\Vec{a_d}(x_c)  
\begin{bmatrix}
|| sin(x_c+ct) || \; cos(c \Delta t) \\
|| cos(x_c+ct) || \; sin(c \Delta t)
\end{bmatrix}
f(x_c).
\end{equation}

Then, working as before, the solution of \ref{lsadv} is given by

\begin{equation}
    \Vec{v}=[{\Psi_1}^T {\Psi_2}^T]
    \begin{bmatrix}
{\sigma_1}^{-1} & 0 \\
0 & {\sigma_2}^{-1}
\end{bmatrix}
\Vec{a_d}(x_c)  
\end{equation}

Substituting the corresponding expressions, we get

\begin{equation}
    v(k)=f(x_c)f(k)\left( \frac{cos(k \Delta x) cos(c \Delta t)}{{|| f(i) cos(i \Delta x) ||}^2 }+\frac{sin(k \Delta x) sin(c \Delta t)}{{|| f(i) sin(i \Delta x) ||}^2 } \right)
\end{equation}


Once again, we consider the simple advection equation solution by setting $f(x)=1$ and a symmetric scheme with $n=1$. Then:

\begin{align}
\label{discsch}
v(1)={}& \frac{cos(\Delta x) cos(c \Delta t)}{2 cos(\Delta x)^2+1}+\frac{sin(c \Delta t)}{2 sin(\Delta x)}\\
v(0)={}& \frac{cos(c \Delta t)}{2 cos(\Delta x)^2+1}\\
v(-1)={}& \frac{cos(\Delta x) cos(c \Delta t)}{2 cos(\Delta x)^2+1}-\frac{sin(c \Delta t)}{2 sin(\Delta x)}    
\end{align}

We can already make several observations at this point. Firstly, it is evident that scheme \ref{discsch} is not symmetric, apart from the limit when $\Delta x \to \pi/2$. Furthermore, $v(0)$ will be equal to $1$ only at the limit of $\Delta x \to \pi/2$ and $c \Delta t \to 0$. Therefore, the least-squares solution will never retrieve the known implicit scheme $[-(c \Delta t)/(2 \Delta x),$ $\; 1, \; (c \Delta t)/(2 \Delta x)]$.

Consistency is also not guaranteed. Requiring that $\sum_i{v(i)}=1$ (for discrete time) and taking the Taylor series expansion of the trigonometric functions, gives a non-trivial relation of $\Delta x$ and $c \Delta t$:

\begin{equation}
    \label{cons}
    \Delta t \approx \frac{2 {\Delta x}^2}{c (3+ \Delta x^2)}
\end{equation}


Finally, when discretizing time and space, a clear bound is observed considering the ability to infer the scheme from the data, given a selected stencil. A data-related CFL condition arises, linked to the structure of the least-square problem. This will be discussed in detail, in the next section.




The observation that scheme \label{discsch} will practically not converge to the known implicit scheme, gives a warning in the case of parametric model inference in discrete time: If the model parametric dependence is assumed to be the one originating from an implicit time discretization scheme, the inference process might not always yield the desired results. In this simple advection case, the inferred operators are not be able to mimic a time-implicit structure of the true discretized operators.


\subsection{Linear Advection Equation}
We try to generalize the above for data that originate from a linear advection problem:

\begin{equation}
    \frac{du}{dt}=c\frac{du}{dx}
\end{equation}

We form a snapshot matrix for $u(x,t)$, over a symmetric neighbourhood of size $(2k+1) \delta x$ of a chosen $x_c$. That is, the snapshot matrix $\mathcal{D}$ is written as:

\begin{equation}
    \mathcal{D}=[u(x_c-k \delta x,t) ... u(x_c,t) ...  u(x_c+k \delta x,t)]
\end{equation}

The 1st SVD mode corresponds to the mean, i.e. $u(x_c,t)$. If additionally, the time interval is selected such that $T=\frac{kL}{c}$, then we can show that $\frac{d^n u}{dt}$ for any $n$ are orthogonal. Since $u(x_c,t)$ is the first SVD mode, then necessarily $\frac{d^n u}{dt}$ correspond to the $n$ other SVD modes.






  
\end{document}
